\documentclass[addpoints]{exam}
% \documentclass[addpoints,12pt]{exam}

\usepackage[slovene]{babel}
\usepackage{amsfonts}
\usepackage{amssymb}
\usepackage{amsmath}
\usepackage{float}
\usepackage{graphicx}
\usepackage{enumitem}
\usepackage{color}
\usepackage{eurosym}

\usepackage{pgfpages}
% \usepackage{tikz}
\usepackage{wrapfig}
\usepackage{pgfkeys}
\usepackage{pgfplots}
\usepackage{xcolor}
\usepackage{tkz-euclide}
\usepackage{xfp}




%Komentar naslednje vrstice glede na verzijo testa -- rešitve ali prostor za odgovore:
% \printanswers

% \linespread{2}


\pointsinrightmargin
\marginbonuspointname{}


\renewcommand{\solutiontitle}{\noindent\textbf{Rešitve in točkovanje:}\par\noindent}

\colorgrids
\definecolor{GridColor}{gray}{0.7}
\setlength{\gridsize}{7mm}

\extrawidth{5mm}
\setlength{\rightpointsmargin}{1.5cm}

\begin{document}

\runningfooter{}{}{}

\begin{center}
    \fbox{\fbox{\parbox{15cm}{\centering
    Preverjanje pred prvim pisnim ocenjevanjem znanja \\ 2. letnik \\ 8. oktober 2025 \\ (Geometrija v ravnini)}}}
\end{center}

\vspace{0.1in}
\makebox[\textwidth]{Ime in priimek, oddelek:\enspace\hrulefill}


\begin{center}
    \hqword{Naloga}
    \hpword{Možne točke}
    % \hbpword{Dodatne točke}
    \hsword{Dosežene točke}
    \htword{Skupaj}  
    % \combinedpointtable[h][questions]
    \gradetable[h][questions]

\end{center}

\vspace{0.1in}


\begin{questions}


    % 1. vprašanje

    \question[4]
    Izračunajte število diagonal in velikost zunanjega kota v pravilnem $12$-kotniku.

    % \begin{solution}[5cm]
        
    % \end{solution}



    % 2. vprašanje

    \question[5]
    Razlika dveh suplementarnih kotov $\rho$ in $\psi$ je $27^\circ 27' 46''$.
    Določite velikosti kotov $\rho$ in $\psi$.

    % \begin{solution}[10cm]
        
    % \end{solution}
    

    % \newpage
    % 3. vprašanje

    \question[6]
    Konstruirajte paralelogram s podatki $b=4~cm$, $v_b=3.5~cm$ in $\alpha=120^\circ$.

    % \begin{solution}[15cm]
        
    % \end{solution}





     % 4. vprašanje

    \question[5]
    Izračunajte velikosti neznanih kotov na skici.

    \begin{figure}[H]
        \centering
        \includegraphics[scale=0.4]{C:/Users/uporabnik/OneDrive/Namizje/SŠ - MAT/GAA/Slike_in_skice/trikotnik_koti.png}
    \end{figure}

    % \begin{solution}[3cm]
        
    % \end{solution}






    % \newpage
    % 5. vprašanje

    \question[3]
    Kolikšna je višina na osnovnico $c$ v enakokrakem trikotniku s krakom $a$ in podatki $c=4\sqrt{2}$ in $a=3\sqrt{3}$?
        % \begin{solution}[3cm]
            
        % \end{solution}



    % 6. vprašanje

    \question[4]
    % Konstruirajte paralelogram s podatki $a=5~cm$, $b=3~cm$, $v=2.5~cm$ ter $\beta>90^\circ$.
    Dan je pravokotni trikotnik s podatki $v=3$ in $a=6$, v katerem leži pravi kot ob oglišču $C$. 
    Izračunajte dolžino druge katete tega trikotnika.
    

    % \begin{solution}[9cm]
        
    % \end{solution}

    % \newpage
    % dodatno vprašanje
    % \vspace{0.1in}
    % \makebox[\textwidth]{Ime in priimek, oddelek:\enspace\hrulefill}

    % ~




\end{questions}



\end{document}